\documentclass[12pt]{article}
\usepackage[english]{babel}
\usepackage[iso,english]{isodate}
\usepackage[table,svgnames]{xcolor}
\usepackage{url}
\usepackage[utf8x]{inputenc}
\usepackage[T1]{fontenc}
\usepackage{longtable}
\usepackage{amsmath}
\usepackage{graphicx}
\usepackage{parskip}
\usepackage{fancyhdr}
\usepackage{vmargin}
\usepackage{hyperref}
\usepackage{pgfgantt}
\usepackage{pgf-umlcd}
\usepackage{xparse}
\usepackage{float}
\usepackage{tabularx}
\usepackage{titling}
\usepackage[nogroupskip,toc,acronym]{glossaries}
\usepackage{fancyhdr}
\usepackage[authoryear]{natbib}
\usepackage{enumitem}
\usepackage{rotating}
%\bibliographystyle{fontysIEEtranN}

\pagestyle{fancy}
\fancyhf{}
\lhead{\includegraphics[width=0.2\linewidth]{TreewatchLogo.pdf}}
\rhead{Personal dossier Jan Kerkenhoff}
\rfoot{Page \thepage}

% alternate row colors for all tables
\definecolor{lightGrey}{rgb}{0.95,0.95,0.95}

\let\oldtable\table
\let\endoldtable\endtable
\renewenvironment{table}{\rowcolors{2}{lightGrey}{}\oldtable}{\endoldtable}

\let\oldtabular\tabular
\let\endoldtabular\endtabular
\renewenvironment{tabular}{\rowcolors{2}{lightGrey}{}\oldtabular}{\endoldtabular}

\let\oldtabularx\tabularx
\let\endoldtabularx\endtabularx
\renewenvironment{tabularx}{\rowcolors{2}{lightGrey}{}\oldtabularx}{\endoldtabularx}

\let\oldlongtable\longtable
\let\endoldlongtable\endlongtable
\renewenvironment{longtable}{\rowcolors{2}{lightGrey}{}\oldlongtable} {\endoldlongtable}

\graphicspath{{../img/}{../img/personaldossiers/}}

\renewcommand{\arraystretch}{1.5}

\DeclareDocumentCommand{\newdualentry}{ O{} O{} m m m m } {
  \newglossaryentry{gls-#3}{name={#5},text={#5\glsadd{#3}},
    description={#6},#1
  }
  \makeglossaries
  \newacronym[see={[Glossary:]{gls-#3}},#2]{#3}{#4}{#5\glsadd{gls-#3}}
}


\newlist{SMART}{description}{2}
\setlist[SMART]{leftmargin=12em,style=nextline}

\newlist{STAR}{description}{2}
\setlist[STAR]{leftmargin=8em,style=nextline}


%\loadglsentries{../glossary.tex}

\setmarginsrb{3 cm}{2.5 cm}{3 cm}{2.5 cm}{1 cm}{1.5 cm}{1 cm}{1.5 cm}

%\makeglossaries
\begin{document}

%%%%%%%%%%%%%%%%%%%%%%%%%%%%%%%%%%%%%%%%%%%%%%%%%%%%%%%%%%%%%%%%%%%%%%%%%%%%%%%%%%%%%%%%% Preface of the report

    \pagenumbering{Roman} % Roman numerals for page counter


    % Title Page
    \begin{titlingpage}
        \begin{center}
            \begin{minipage}{\linewidth}
            \centering
            %University logo
            \includegraphics[width=0.3\linewidth]{FontysLogo.pdf}
            \par
            \vspace{3cm}
            %Thesis title
            {\uppercase
                {\Large Personal dossier\\ Jan Kerkenhoff \\ 2015 Sofa GTL \\ TreeWatch Project
            \par
            \vspace{3cm}}}
            \includegraphics[width=0.5\linewidth]{TreewatchLogo.pdf}
            \par
            \vspace{2cm}
            %Author's name
            {Jan Kerkenhoff\\
	         \par}
            \vspace{2cm}

            %Date
            \today
            \end{minipage}
        \end{center}
    \end{titlingpage}
    \clearpage

\setmarginsrb{3 cm}{2.5 cm}{3 cm}{2.5 cm}{1.5 cm}{1.5 cm}{1 cm}{1.5 cm}
    \section*{Document information}
\addcontentsline{toc}{section}{Document information}
\renewenvironment{tabular}{\oldtabular}{\endoldtabular}
	\begin{tabular}{ll}
		\textbf{Document name:} & Personal dossier Jan Kerkenhoff\\
		\textbf{Document owner:} & Jan Kerkenhoff \\
		\textbf{Company/Organisation:} & Fleuren Baarlo \\
		\textbf{Contact person:} & Jan Kerkenhoff \\
		\textbf{Date:} & \today \\
		\textbf{Place:} & Fontys University of Applied Science Venlo \\
		\textbf{Author:} & \parbox[t]{5cm}{
		Jan Kerkenhoff\\ j.kerkenhoff@student.fontys.nl\\ 2197164 \\}
	\end{tabular}
\renewenvironment{tabular}{\rowcolors{2}{lightGrey}{}\oldtabular}{\endoldtabular}

    %\printglossary[type=\acronymtype]
    %\printglossary
    \pagebreak

    %\listoffigures
    %\addcontentsline{toc}{section}{\listfigurename}
    %\listoftables
    %\addcontentsline{toc}{section}{\listtablename}
    %\pagebreak

    \tableofcontents
    \clearpage

%%%%%%%%%%%%%%%%%%%%%%%%%%%%%%%%%%%%%%%%%%%%%%%%%%%%%%%%%%%%%%%%%%%%%%%%%%%%%%%%%%%%%%%%% Main Part of the Report

    \pagenumbering{arabic}
   
    \begin{sidewaysfigure}[ht]

  	    \centering
  	    \includegraphics[width=\textwidth, keepaspectratio=true]{LearningGoalProgressJK.png}
  	    \caption{Learning goal progress}\label{lgp}
    \end{sidewaysfigure}
    
    \section{A-Competencies}
    
    \subsection{SMART Learning Goal 1}
\begin{SMART}
    \item[Specific] As the head developer for this Sofa I want to get more experience developing applications for mobile systems like Android and iOS. specializing in a framework for running web apps native on different mobile systems
    \item[Measurable] To measure this learning goal I will take into account the number of errors per lines of code.
    \item[Attainable] The topic is widely described in the internet, there are many useful tutorials and guides available, also an extensive documentation is provide for almost all technologies
    \item[Relevant] In my future career, in which I want to become a application developer mobile development is often demanded skill.
    \item[Time-limited] By January 2016, this project will be closed, Therefore I will be able to develop mobile for the project until then.
    \item[My complete goal] To improve my mobile development experience by January 2016.
\end{SMART}

\subsection{STAR Analysis 1}
\begin{STAR}
    \item[Situation] We used Xamarin and the Xamarin Forms SDK to create our cross platform application. I was responsible for coordinating the development across the team, also researched and picked out the dependencies we used. The accompanying learning goal is learning goal 1.
    \item[Task] I wanted to improve my experience with developing mobile apps. Especially with Xamarin Forms and Xamarin in general, since it was used in the project.
    \item[Action] I worked on several different parts of the application.
    The first part was showing a map on iOS and Android. This was accomplished by writing a custom renderer as seen in CODEREF-MAPRENDERER. This render used the platform native app implementations to show the map.
    We later added the functionality to add overlays of the field to the map and last but not least we added heat map overlays too. This is also done in the aforementioned renderer, specific for each platform. The heat map part was only done on iOS, because we encountered severe performance issues. To combat these issues I created a new Polygon class on iOS that can draw multiple polygons at once, this can be see in CODEREF-MULTIPOLYGON. 
    I also worked together with René Karoff on the KML part of our application. This was done to load real data from our customer in the app. For this we created a kml-parser as seen in CODEREF-KMLParser. This allowed us to read in real data about the field from fleuren, like the different blocks, the field boundaries or heat maps.
     
    \item[Result] I learned a lot about developing mobile apps with Xamarin Forms, also gained additional knowledge with C\# which is used by Xamarin. As can seen in figure~\ref{lgp} the measurement for the accompanying learning goal, the errors per 50 Lines of code, also declined through out the project.
\end{STAR}

\subsection{SMART Learning Goal 2}
\begin{SMART}
    \item[Specific] I want to get more experience in doing test driven development. This project is involving different technologies for mobile development and web development, I want to learn more about testing for the web and mobile devices. Including but not limited to regression testing and maybe some kind of fuzzing.
    \item[Measurable] To measure this learning goal I will take into account the number of tests if have written and executed this number will be delivered by the number of tests that actually break the software.
    \item[Attainable] The topic is widely described in the internet, there are many useful tutorials and guides available, also an extensive documentation is available. Some technologies are quite new and not that much is information available but there still is enough to get started
    \item[Relevant] Test driven development is the way to go in the industry when I comes to develop new applications.
    \item[Time-limited] By January 2016, this project will be closed, Therefore I will be able to manage testing for the project until then.
    \item[My complete goal] To improve my testing  experience by January 2016.
\end{SMART}

\subsection{STAR Analysis 2}
\begin{STAR}
    \item[Situation] We decided to use Test Driven Development as part of our quality management. 
    \item[Task] I wanted to gain more knowledge about TDD with mobile and about TDD in general. 
    \item[Action] I researched a lot about the technologies that could have been used for TDD with Xamarin. I always worked together closely with our quality manager in this tasks. We use normal unit tests via the NUnit library for everything that is non GUI. For functional GUI tests we choose xamarin own UiTest library. This allowed us to test both backed and fronted of our application. I have written several unit tests for our utility classes like the one for the CODEREF-UNITTESTS. I also did our first batch of UiTests for the basic functions at the start of the project. You seen an example of UiTests in CODEREF-UITests I refined these throughout the project together with the other team members to test additional GUI elements. I also setup our Continues integration environment with AppVeyor, which always compiled our project and executed the unittests. UiTests were not support by AppVeyor and we executed them manually locally, since we did not have subscription to the Xamarin Test Cloud.
    \item[Result] I learned a lot about the NUnit and UiTests libraries. 
    As can seen in figure~\ref{lgp} the measurements for the accompanying learning goal increased over time. If have written more tests and some even detected hidden flaws in the application we didn't notice before this is counted in "Tests breaking the app". I also learned that it is hard to motivate people to do TDD. A lot of my group members tended to forget about writing tests, I even sometimes forgot it myself while fixing bugs.
    \item[Reflection] I learned a lot about TDD development, but still failed to get the group to continuously use TDD for every step. This is also partly related to the fact that we had a hard time figuring out who to even test certain parts of the app.
\end{STAR}

\section{B-Competencies}

\subsection{SMART Learning Goal 3}
\begin{SMART}
    \item[Specific] I want to improve my skills in presenting my ideas to people.
    \item[Measurable] To measure this learning goal I will take into account the number of different documents, meetings and presentations done for the project.
    \item[Attainable] I already learned a lot in previous projects and the internship. I will try to use the gained knowledge and use that to improve my presentations.
    \item[Relevant] Expressing your thoughts and ideas to other people is always needed in software development.
    \item[Time-limited] By January 2016, this project will be closed, Therefore I will be able improve my communication until then.
    \item[My complete goal] To improve the way I express ideas and thoughts by January 2016.
\end{SMART}

\subsection{STAR Analysis 3}
\begin{STAR}
    \item[Situation] Being in an project with several members requires you to present your ideas to other people. This needs to be done in different ways, one being reports and presentation and also verbally when talking about ideas in the group.
    \item[Task]
    I wanted to improve my communication skills when it comes to presenting my ideas to others. I always had minor problems getting my point across clearly and precisely.
    \item[Action] To achieve this goal I tried to be involved in as many communications and meetings as possible. I used the meetings we had with our customer to talk about the problems we encountered and solutions we picked out. I also took part in creating our presentations and documents. I wanted to get as much chances as possible to improve myself.
    \item[Result]
    To track the progress of this goal I used a short questionnaire that asses how good my other team members understood my ideas when I presented them to them. This is can be tracked in figure~\ref{lgp}, I used a scale from 1 being "only understand parts of the idea" to 10 being "I understand the whole idea and know exactly what you want to". I asked my group members to fill out this survey several time throughout the project to measure my progress. As seen in figure~\ref{lgp} my communication skills went up and is was able to better bring my point across towards the end of the project.
    \item[Reflection] I improved my communication quite good. I did get some chances to talk to our customer which was really nice. I feel pretty good about my english right now, using it in a work environment and so on.
\end{STAR}


%\clearpage

%\bibliography{biblist}

\end{document}
